\documentclass[12pt, a4paper]{article}

% ============================================
% PACKAGE IMPORTS
% ============================================

% Encoding and language
\usepackage[utf8]{inputenc}
\usepackage[T1]{fontenc}
\usepackage[english]{babel}

% Page layout
\usepackage{geometry}
\geometry{margin=1in}

% Typography
\usepackage{setspace}
\onehalfspacing

% Graphics and figures
\usepackage{graphicx}
\usepackage{float}
\graphicspath{{images/}}

% Tables
\usepackage{booktabs}
\usepackage{multirow}

% Math
\usepackage{amsmath, amssymb, amsthm}

% Colors
\usepackage{xcolor}

% Hyperlinks
\usepackage{hyperref}
\hypersetup{
    colorlinks=true,
    linkcolor=blue,
    filecolor=magenta,
    urlcolor=cyan,
    citecolor=blue
}

% Headers and footers
\usepackage{fancyhdr}
\pagestyle{fancy}
\fancyhf{}
\rhead{\thepage}
\lhead{\leftmark}

% Bibliography
\usepackage[style=numeric, sorting=none]{biblatex}
\addbibresource{references.bib}

% ============================================
% DOCUMENT INFORMATION
% ============================================

\title{Your Report Title Here}
\author{Your Name\\
        \texttt{your.email@example.com}}
\date{\today}

% ============================================
% DOCUMENT CONTENT
% ============================================

\begin{document}

% Title page
\maketitle
\thispagestyle{empty}

% Abstract
\begin{abstract}
    This is the abstract of your report. It should provide a brief overview of the main points, 
    findings, and conclusions of your work. Keep it concise and informative, typically 
    150-250 words.
\end{abstract}

\newpage

% Table of contents
\tableofcontents
\newpage

% ============================================
% MAIN CONTENT
% ============================================

\section{Introduction}
\label{sec:introduction}

This is the introduction section. Here you should:
\begin{itemize}
    \item Introduce the topic and provide background information
    \item State the purpose and objectives of the report
    \item Outline the structure of the document
\end{itemize}

You can reference other sections like Section~\ref{sec:methodology}.

\section{Methodology}
\label{sec:methodology}

Describe the methods and approaches used in your work. This section might include:

\subsection{Research Design}
Explain the overall research design and approach.

\subsection{Data Collection}
Detail how data was collected or what sources were used.

\subsection{Analysis Methods}
Describe the analytical methods or tools employed.

\section{Results}
\label{sec:results}

Present your findings here. You can include:

\subsection{Tables}

Example table:
\begin{table}[H]
    \centering
    \begin{tabular}{lcc}
        \toprule
        \textbf{Item} & \textbf{Value 1} & \textbf{Value 2} \\
        \midrule
        Sample A & 10.5 & 12.3 \\
        Sample B & 8.7 & 9.1 \\
        Sample C & 15.2 & 14.8 \\
        \bottomrule
    \end{tabular}
    \caption{Example table showing sample data.}
    \label{tab:example}
\end{table}

\subsection{Figures}

Example figure reference (uncomment when you have images):
% \begin{figure}[H]
%     \centering
%     \includegraphics[width=0.7\textwidth]{example-image}
%     \caption{Example figure caption.}
%     \label{fig:example}
% \end{figure}

\subsection{Mathematical Equations}

You can include equations:
\begin{equation}
    E = mc^2
    \label{eq:einstein}
\end{equation}

Where $E$ is energy, $m$ is mass, and $c$ is the speed of light.

\section{Discussion}
\label{sec:discussion}

Interpret and discuss your results. Consider:
\begin{itemize}
    \item What do the results mean?
    \item How do they relate to existing knowledge?
    \item What are the implications?
    \item What are the limitations?
\end{itemize}

\section{Conclusion}
\label{sec:conclusion}

Summarize the main findings and their significance. Include:
\begin{itemize}
    \item Brief recap of objectives
    \item Summary of key findings
    \item Final thoughts and recommendations
    \item Suggestions for future work
\end{itemize}

% ============================================
% BIBLIOGRAPHY
% ============================================

\newpage
\printbibliography[heading=bibintoc, title={References}]

% ============================================
% APPENDICES (Optional)
% ============================================

\newpage
\appendix

\section{Additional Data}
\label{app:additional-data}

Any supplementary information, additional tables, or detailed calculations can go here.

\section{Code Listings}
\label{app:code}

If you need to include code or technical details, place them in appendices.

\end{document}
